\chapter{Proposed Hybrid Architecture}
\label{ch:architecture}

\section{From received observations to an alarm}
The detector has three branches: General, Drift and Noise. General evaluates broad deviations using a mixture of residual classifiers. Drift and Noise combine causal temporal evidence with observations collected during a declared three-hour decision interval. A network-level evidence router, local feedback heads and a verifier regulate the specialist alarms. The final prediction is the union of General's alarm and the specialist alarms that survive these checks.

Feedback and verification can reject specialist additions, but every General alarm remains in the final output. The experiments assess whether restricting those additions improves the balance of detections and errors.

Figure~\ref{fig:finalflow} shows the information paths and the two routers. General's internal router weights its five component predictions at one observed endpoint. The external evidence router describes the likely mechanism at a network-time group and informs specialist rejection. The early-warning head is another network-time predictor, but uses causal inputs only. Its outputs are inputs to the verifier, not a separate branch of sensor alarms.

\figfile{final_flow.pdf}{1.0}{Detailed information paths in the hybrid architecture}{Information flow in the selected architecture. Arrows show information dependencies. Only causal specialist scores enter the early-warning head. General remains available independently of those checks. The received-history extension enters General in L-Town. Both networks retain the three-branch organization.}{fig:finalflow}

For an observed pressure endpoint $(i,t)$, let $x_{i,t}$ denote the common causal feature vector. Source and scenario identifiers separate trajectories during feature construction and are excluded from the classifier inputs. Training labels and event identities determine supervision and sample weights. At inference, the prediction path consumes received observations, masks, time alignment and fitted parameters. The source-code and configuration map in Appendix~\ref{app:evidence} identifies the implementation behind the equations in this chapter.

\section{A target-group-excluding normal reference}
\subsection{Fitting the normal representation}
The reference estimates a pressure reading from other observed sensors. It is fitted on normal TRAIN observations, including their measurement noise and missingness. For sensor $i$, the base fit stores an observed-data mean $\mu_i$ and a scale $a_i$ representing its normal variation. It standardizes available values as $(y_{i,t}-\mu_i)/a_i$ and performs iterative low-rank completion. At each iteration, the observed entries remain fixed and only missing entries are replaced by the low-rank reconstruction.

The retained basis $B\in\mathbb R^{N\times16}$ represents normal cross-sensor variation. Twelve completion iterations are used. Sensors are assigned to four groups by their fixed index modulo four. When the reference predicts a group, every member of that group is excluded from the supporting observations. This also excludes the target reading from the robust reweighting calculations used for its own prediction.

The base reference first estimates a prediction-error scale $\sigma_i^{(0)}$ from normal observed residuals. The calculation uses the median absolute prediction error divided by $0.67448975$, with a positive floor. This quantity scales prediction discrepancies; it differs from the wider hydraulic-variation scale $a_i$. The robust wrapper subsequently recalibrates a final error scale $\sigma_i$ on normal training observations. These fitted quantities remain fixed during calibration and evaluation.

\subsection{Robust inference from visible sensors}
For a target group $g$, the robust reference uses three fixed views of its permitted supporting sensors. One includes every sensor outside $g$; two retain random subsets with an 80\% inclusion probability, generated with seed 821. A subset with insufficient support for the rank falls back to the full permitted set. The same view masks are reused across prediction times.

Equation~\eqref{eq:whitenedreference} defines the whitened design row and observed response:
\begin{equation}
 d_j=\frac{a_j B_{j,:}}{\sigma_j^{(0)}},
 \qquad r_{j,t}=\frac{y_{j,t}-\mu_j}{\sigma_j^{(0)}}.\label{eq:whitenedreference}
\end{equation}
For view $v$, the latent vector solves the weighted ridge problem in Equation~\eqref{eq:referencefit}:
\begin{equation}
 \widehat z_{g,v,t}=\arg\min_z
 \left\{\sum_{j\notin g} w_{j,v,t}(r_{j,t}-d_jz)^2
       +\lambda\lVert z\rVert_2^2\right\},\qquad\lambda=10^{-3}.
 \label{eq:referencefit}
\end{equation}
A missing or excluded observation has zero weight. Each view performs six solves. Between solves, the implementation reduces the weight of observations with large leverage-adjusted discrepancies, initially with a bounded inverse-error weight and subsequently with a bounded biweight rule. The later rule uses cutoff 4.0 and retains a minimum relative weight of 0.01 for visible observations. Group exclusion remains in force throughout.

Each view gives an estimate $\mu_i+a_iB_{i,:}\widehat z_{g,v,t}$ for a target sensor. The final reference takes the median across the three estimates. The median absolute disagreement among these estimates provides a reliability feature. A separate support feature records the proportion of sensors outside the target group that are currently observed. Support is measured before robust downweighting. Figure~\ref{fig:referenceviews} illustrates target exclusion and the three reference views.

The resulting normalized residual is
\begin{equation}
 u_{i,t}=\frac{y_{i,t}-\widehat p_{i,t}}{\sigma_i}.\label{eq:normalizedresidual}
\end{equation}
Equation~\eqref{eq:normalizedresidual} measures the difference between the observation and its reference in units of the fitted error scale. Robust reweighting reduces the contribution of inconsistent supporting readings. The disagreement and support channels record variation across reference views and the availability of supporting readings.

\figfile{reference_views.pdf}{0.95}{Target exclusion and robust reference views}{Schematic of reference prediction for one target group. Every view excludes the complete target group. Missingness reduces the visible support within each view, and the median combines the three predictions. The group and view masks are illustrative; the numerical fit uses the frozen masks and rank specified in the text.}{fig:referenceviews}

\section{The common evidence bank}
The shared causal bank contains 116 columns. Its first 109 combine residual summaries, adaptive dynamic evidence, reference disagreement and sequential mechanism features. Seven daily comparison features complete the bank. The construction retains observed endpoints from hour 15 onward, providing the initial history required by the base residual window. An absent pressure reading produces no evaluated endpoint at that hour.

Current residual magnitude captures an immediate discrepancy. Windowed means, variation, slopes and sign consistency describe how that discrepancy behaves over time. Dynamic and sequential channels summarize innovations, accumulated directional evidence and excess energy. They let a classifier distinguish a large isolated deviation from a sustained pattern and interpret either in relation to observation support. Appendix~\ref{app:evidence} links to the ordered feature manifests and construction modules in the repository.

The daily comparisons use actually received pressures at lags 24 and 48 hours. In this bank their normalization is based on the declared observation-noise scale $\sigma_p=0.1$ m:
\begin{equation}
 d^{\mathrm{seasonal}}_{i,t,\ell}=
 \frac{y_{i,t}-y_{i,t-\ell}}{\sqrt{2}\sigma_p},
 \qquad \ell\in\{24,48\}.\label{eq:seasonaldifference}
\end{equation}
Each lag in Equation~\eqref{eq:seasonaldifference} supplies a signed difference, its magnitude and a support indicator. If the historical reading is absent, the difference is filled with zero and support is zero. The seventh feature combines the signs and magnitudes of the two supported comparisons:
\begin{equation}
 c^{\mathrm{seasonal}}_{i,t}=\operatorname{sign}(d_{24}d_{48})
                   \min(|d_{24}|,|d_{48}|).\label{eq:seasonalconsistency}
\end{equation}
The signed differences and their combination in Equation~\eqref{eq:seasonalconsistency} exploit the daily context of the generated demand process. Their implementation requires a contiguous hourly scenario grid and checks the alignment of the flattened observed endpoints.

The final General mixture uses the 116-column bank, with feature subsets selected internally. The L-Town history extension increases its input to 158 columns. The specialist causal bank remains at 116 in both networks. Table~\ref{tab:dimensions} lists the feature dimensions at each interface.

\begin{table}[tbp]
\centering\small
\caption[Feature dimensions at each architectural interface]{Dimensions of the selected feature interfaces. A component's internal feature profile may use a subset of its available bank. Graph-time features summarize the observed sensors at one time.}
\label{tab:dimensions}
\begin{tabularx}{\linewidth}{@{}YrY@{}}
\toprule
Interface & Columns & Construction\\
\midrule
Common causal bank & 116 & 109 residual/dynamic/sequential columns plus 7 seasonal columns\\
L-Town General bank & 158 & Common bank plus 42 received-history columns\\
Delayed specialist input & 175 & Common bank plus 48 offset values and 11 summaries\\
External router input & 135 & Mean, standard deviation and maximum of 45 local channels\\
Local feedback input & 50 & 45 local channels plus 5 external-router probabilities\\
Early-warning input & 160 & Causal graph-time aggregates, sensor count and past differences\\
Verifier input & 57 & 37 raw channels, 8 score channels and 12 early-warning channels\\
\bottomrule
\end{tabularx}
\end{table}

\section{General and its five internal components}
\subsection{Component roles and training populations}
General contains five histogram gradient-boosted classifiers: general, abrupt, replay, drift and noise. The first uses all available columns and learns binary corruption detection across every attack family. The other four use mechanism-oriented feature profiles. Abrupt shares current residual and short-window evidence across random and targeted attacks. Replay includes lag-related and longer-window evidence. Drift emphasizes coverage, means, slopes and sign consistency; noise emphasizes coverage and dispersion. Common residual, support and gap features are available to all profiles.

The fit retains every positive TRAIN endpoint and samples at most 60,000 negative endpoints. A sampled negative receives a population-correction weight equal to the total number of negatives divided by the sampled count. The general component preserves this weighted prevalence. Each internal specialist trains on its own mechanism's positives and the sampled negatives, with the two binary classes rebalanced. Other mechanisms' positive endpoints are excluded from that specialist's binary fit.

The internal router uses five classes with a different interpretation from five attack labels. Class zero covers unmodified endpoints; random and targeted positives share the abrupt class, and replay, drift and noise positives occupy the remaining classes. The router is trained on the same sampled population with mechanism-balanced weights. Its class-zero probability weights the general component, which can still assign a high anomaly score. The router probability determines a component's weight; the component score estimates corruption.

All five components and the internal router use the same fixed estimator recipe: 120 boosting iterations, at most 15 leaves, a minimum of 30 samples per leaf, learning rate 0.08 and L2 regularization 10. Early stopping is disabled. The model seed changes the stochastic realization and sample selection; calibration chooses the operating point after fitting.

\subsection{Internal mixture and symmetric calibration}
Let $s_k(x)$ be the five component scores and $\pi_k(x)$ the internal-router probabilities. Equation~\eqref{eq:generalbase} defines the basic General score:
\begin{equation}
 g(x)=\sum_{k=0}^{4}\pi_k(x)s_k(x).\label{eq:generalbase}
\end{equation}
Modena retains this direct mixture in its selected Stage-E system. L-Town uses the saved symmetric router transformation. With temperature $T>0$ and uniform shrinkage $\rho\in[0,1]$,
\begin{align}
 \widetilde\pi_k(x)&=\frac{\max(\pi_k(x),10^{-12})^{1/T}}
 {\sum_j\max(\pi_j(x),10^{-12})^{1/T}},\label{eq:generaltemperature}\\
 w_k(x)&=(1-\rho)\widetilde\pi_k(x)+\rho/5,\label{eq:generalshrinkage}\\
 g(x)&=\sum_{k=0}^{4}w_k(x)s_k(x).
 \label{eq:generalcal}
\end{align}
Temperature in Equation~\eqref{eq:generaltemperature} controls concentration, and shrinkage in Equation~\eqref{eq:generalshrinkage} retains a uniform contribution from all components in the score of Equation~\eqref{eq:generalcal}. The same transformation applies to every class. The selected values and General threshold are stored for each fitted realization. No family label is supplied to this calculation at inference.

\subsection{The received-history extension in L-Town}
The extension adds the same 42 columns to every General component's profile and to its internal router. Equation~\eqref{eq:historylags} gives the lag set:
\begin{equation}
 \mathcal L=\{1,2,\ldots,12,24,48\}\ \mathrm{hours}.\label{eq:historylags}
\end{equation}
For an available historical reading, Equation~\eqref{eq:receivedhistory} defines the quantized difference:
\begin{equation}
 h_{i,t,\ell}=0.1\operatorname{round}
 \left(\frac{y_{i,t}-y_{i,t-\ell}}{0.1\sigma_i}\right).
 \label{eq:receivedhistory}
\end{equation}
Each lag contributes $h_{i,t,\ell}$, $|h_{i,t,\ell}|$ and an availability indicator. Missing historical differences remain missing numerical values, with the indicator set to zero. This encoding differs from the seasonal bank's zero-plus-support convention and uses the fitted reference-error scale rather than $\sqrt{2}\sigma_p$.

Matching uses the exact historical hour within the same scenario and sensor. The retained history is the received stream, so previous corruption remains part of what the detector knows. Quantization removes access to arbitrarily precise equality through these new columns. The extension changes the learned representation inside General; the external specialist pipeline is preserved. The development and confirmation evidence for selecting it is documented in Chapter~\ref{ch:development} and Appendix~\ref{app:evidence}.

\section{Drift and Noise: causal scores and delayed evidence}
\subsection{Complementary causal components}
The Drift branch combines a full seasonal classifier with a tuned residual classifier. The Noise branch combines a fast seasonal classifier with a full seasonal classifier. The fast profile emphasizes recent changes and energy rather than long-lived accumulated states, helping its score return toward normal after a disturbance ends.

The combination operates in log-odds space. Define $L(p)=\log(p/(1-p))$, with probabilities clipped to $[10^{-6},1-10^{-6}]$, and let $\operatorname{sigmoid}$ be its inverse. The causal branch scores are
\begin{align}
 c_D&=\operatorname{sigmoid}\left(0.90L(s_{D,\mathrm{seasonal}})
                                  +0.10L(s_{D,\mathrm{tuned}})\right),\label{eq:causaldrift}\\
 c_N&=\operatorname{sigmoid}\left(0.95L(s_{N,\mathrm{fast}})
                                  +0.05L(s_{N,\mathrm{full}})\right).\label{eq:causalnoise}
\end{align}
The implementation clips the combined logits in Equations~\eqref{eq:causaldrift} and~\eqref{eq:causalnoise} to $[-30,30]$. The coefficients are part of the frozen recipe, while the component classifiers are refitted for each model seed.

The seasonal classifiers use LightGBM. Full models have 400 trees and at most 31 leaves; the fast noise model has 350 trees and at most 23 leaves. Their training target is an altered endpoint of the assigned family. Training weight is divided between own-family positives (0.45), unmodified endpoints in that family's events (0.30), clean endpoints (0.15) and other-family endpoints (0.10). Weights are further balanced across events or scenarios within each category. Early positive endpoints receive a multiplier of two before within-event normalization. This weighting trains the specialist to detect its assigned family while rejecting competing mechanisms.

\subsection{A bounded observation interval}
For endpoint $(i,t)$, the delayed feature constructor gathers 16 specified causal columns at offsets $t+1$, $t+2$ and $t+3$. These include residuals, rates, support, innovations, short-window drift/noise evidence and seasonal comparisons. Exact offset matching yields 48 values. Eleven additional summaries describe the available residual trajectory over $[t,t+3]$: extrema, mean, dispersion, rise, slope, absolute-value summaries, support and sign agreement.

A missing offset stays missing. Summary statistics use available values, and slope uses their actual hour offsets. If there is no temporal variation in the available offsets, slope remains unavailable. The current observed endpoint always anchors the window. Near the end of a finite scenario, the operator uses only the permitted observations present in that scenario; it never crosses into another trajectory.

Separate delayed Drift and Noise classifiers consume the 175-column augmented bank. Each uses the full seasonal LightGBM recipe and family-weighting procedure. In parallel, a forward maximum pools the causal branch scores over available endpoints in the same bounded interval:
\begin{equation}
 m_b(i,t)=\max_{u\in[t,t+3]:\,(i,u)\ \mathrm{available}}c_b(i,u),
 \qquad b\in\{D,N\}.\label{eq:forwardmaximum}
\end{equation}
The final specialist score combines the learned delayed score $d_b$ and this maximum with equal logit weights:
\begin{equation}
 q_b(i,t)=\operatorname{sigmoid}\left(\tfrac12L(d_b(i,t))+\tfrac12L(m_b(i,t))\right).
 \label{eq:specialistfinal}
\end{equation}
The delayed classifier summarizes the local trajectory, while Equation~\eqref{eq:forwardmaximum} retains the strongest causal score within the allowance. Equation~\eqref{eq:specialistfinal} combines these two scores. Both finalize the decision about $t$ on the $t+3$ clock.

\section{External routing and local feedback}
The external router receives 37 designated raw evidence channels and eight specialist score channels: causal, forward-maximum, delayed and final scores for each branch. For each source/scenario/time group, the 45 local channels are aggregated by mean, population standard deviation and maximum across the retained observed endpoints. The resulting 135-dimensional vector produces five probabilities for clean, abrupt, replay, drift and noise evidence.

Because these inputs include delayed scores, this router runs on the specialized decision clock. Its probabilities are broadcast to the observed endpoints in that group. Two feedback classifiers then combine the 45 local channels with the five router probabilities. Each estimates whether the endpoint belongs to its own specialist's positive target. Training includes own-family positives, positives from other mechanisms and sampled negatives, including high-scoring negatives. It uses source-held router probabilities to construct the feedback training inputs.

The selected veto rejects specialist evidence when replay is more likely and local feedback does not support the specialist mechanism. Let $r_R$, $r_D$ and $r_N$ denote the external probabilities for replay, drift and noise, and let $f_b$ denote local feedback. Then
\begin{equation}
 V_b(i,t)=\mathbf1\{r_R>r_b+0.10\}\,
          \mathbf1\{f_b(i,t)<0.10\},\qquad b\in\{D,N\}.
 \label{eq:veto}
\end{equation}
Both conditions in Equation~\eqref{eq:veto} must hold: the network context must favor replay and local feedback must reject the specialist mechanism. The guard addresses overlap between replay and specialist evidence. The received-history extension retains this rule.

The veto can remove a specialist alarm but cannot generate an alarm or cancel General. Feedback is a learned check on expert evidence, with parameters fixed during evaluation.

\section{Causal early warning and specialist verification}
\subsection{A separate network-time signal}
The early-warning head uses the same 37 raw evidence channels and only the two causal specialist scores. Their three graph-time aggregates give 117 values. The observed-sensor count adds one, and differences in 13 selected mean channels at past lags 1, 2 and 3, with three availability indicators, complete the 160-column input.

The head predicts the five evidence classes using a histogram gradient-boosted classifier. A confidence threshold permits abstention. In the training procedure, the first three hours of an event receive extra weight, and source-held predictions determine the abstention setting under a clean-warning budget. Event onset is a training/evaluation annotation; prediction uses the causal feature vector.

A warning describes the network-time group; it does not identify which sensor is corrupted. For final sensor decisions, warning probabilities, confidence, abstention, persistence and one- and two-hour Drift/Noise probability lags become twelve verifier inputs. Lack of a warning is not implemented as a mandatory rejection rule, although the learned verifier can use the warning features when judging a specialist candidate.

\subsection{Verification on the candidate region}
Each verifier head receives 57 columns: 37 raw channels, eight specialist scores and twelve early-warning channels. Its training population is the high-score candidate region defined from clean TRAIN out-of-fold scores. The pre-threshold corresponds to the upper two-percent region of that branch's clean score distribution. This selects the region in which specialist precision matters most.

The target is an endpoint modified by an attack from the verifier's assigned family. Other-family positives therefore serve as negatives for that branch-specific verification task, as do unmodified endpoints. A class-balanced histogram boosting model with 200 iterations and at most 23 leaves fits each head. The region boundary and fitted model are then frozen.

Let $v_b(x)$ be the verifier's keep probability within its region. Outside that region the implementation returns one, leaving those scores unaffected. At a calibrated cutoff $\kappa_b$, the verification mask is
\begin{equation}
 K_b(i,t)=\mathbf1\{v_b(i,t)\geq\kappa_b\}.\label{eq:verifiermask}
\end{equation}
A cutoff of zero in Equation~\eqref{eq:verifiermask} disables the check for that branch. General is never passed through a verifier. The selected cutoff therefore determines whether each fitted verifier head is active.

\section{The final decision and its calibration}
Given fitted scores and frozen thresholds $\tau_G,\tau_D,\tau_N$, define
\begin{align}
 A_G(i,t)&=\mathbf1\{g(i,t)>\tau_G\},\label{eq:generalalarm}\\
 A_b(i,t)&=\mathbf1\{q_b(i,t)>\tau_b\}\,(1-V_b(i,t))K_b(i,t),\quad b\in\{D,N\},\label{eq:specialistalarm}\\
 \widehat y_{i,t}&=A_G(i,t)\lor A_D(i,t)\lor A_N(i,t).
 \label{eq:finaldecision}
\end{align}
The immediate path in Equation~\eqref{eq:generalalarm} can emit $A_G$ at $t$. Specialist acceptance in Equation~\eqref{eq:specialistalarm} can add an alarm for that endpoint after the declared interval; Equation~\eqref{eq:finaldecision} takes their union. Relative to the same General score and threshold, the final positive set always contains General's positive set. Guards and verification only restrict the specialist additions. Changes to General's representation or threshold require a separate paired comparison because they change this reference set. Table~\ref{tab:algorithm} summarizes the endpoint calculation.

\begin{table}[!htbp]
\centering\small
\caption[Computation of one endpoint decision]{Logical evaluation for endpoint $(i,t)$. History/state construction is shared across endpoints; the steps specify information dependencies rather than a required software execution order.}
\label{tab:algorithm}
\begin{tabularx}{\linewidth}{@{}rY@{}}
\toprule
Step & Operation\\
\midrule
1 & Build the target-excluding reference and the causal feature vector from received observations and frozen parameters.\\
2 & Compute the five General component scores, internal weights and immediate alarm $A_G$. Include the 42 received-history columns for L-Town.\\
3 & Retain causal specialist and early-warning outputs; collect permitted specialist evidence through $t+3$.\\
4 & Form forward maxima, delayed scores and final specialist scores $q_D,q_N$.\\
5 & Evaluate external routing, local feedback and the two verifier probabilities. Apply each veto and verifier once.\\
6 & Return the union in Equation~\eqref{eq:finaldecision}, retaining the endpoint timestamp and declared decision time.\\
\bottomrule
\end{tabularx}
\end{table}

Threshold selection is performed on calibration. The Stage-E search allocates a clean false-alarm budget of 0.005 across the three raw branch thresholds, then evaluates their guarded union. The selected candidate maximizes pooled F1 subject to the recorded family protections, with worst-family F1 used in its ordering of tied candidates. Verifier cutoffs are chosen on the declared grid. Calibration may use less than this budget. The ceiling applies to calibration FPR and does not guarantee the evaluation rate.

L-Town's protected-history selection holds the specialized alarm path fixed while choosing General's symmetric router settings and threshold. Its protection checks compare the complete candidate with its paired baseline, including source-specific checks. The resulting operating points are retained separately for each model seed. Appendix~\ref{app:evidence} gives an example and links to the complete configuration records in the repository.

Source-held fitting supports the integration layers. Each TRAIN source is held out in turn; the reference and base experts used for its out-of-fold predictions are fitted without that source. These out-of-fold scores provide training inputs for the warning, feedback and verifier stages. Delayed classifiers are fitted within each fold alongside the base experts. Final classifiers are then fitted for the selected recipe, and calibration selects decision settings. The normal reference and the deterministic feature caches are shared across model-seed repetitions, while the stochastic classifiers and their relevant samples are refitted.

\section{The two selected implementations}
Modena retains the confirmed Stage-E candidate: the 116-column General mixture, the two specialized branches and the guarded decision layer with early-warning-assisted verification. L-Town retains the confirmed full-history General with 158 available columns, symmetric internal-router calibration and the preserved specialized pipeline. Both use the same functional separation between broad detection, temporal specialization and controlled admission of specialist alarms.
